\documentclass[a4paper]{article}     
\usepackage{amsfonts}
\usepackage{amsmath}
\usepackage{amssymb}
\usepackage{graphicx}
\usepackage{cancel}

\newcommand{\asa}{\Leftrightarrow} %als slechts als
\newcommand{\adR}{\Rightarrow} %als dan Rechts
\newcommand{\adL}{\Leftarrow}
\newcommand{\Lh}{\left(} %Links haakje
\newcommand{\Rh}{\right)}
\newcommand{\Lv}{\left[} %Links vierkanthaakje
\newcommand{\Rv}{\right]}
\newcommand{\La}{\left\{} %Links acolade
\newcommand{\Ra}{\right\}}
\newcommand{\Ras}{\right.} %Rechts acolade stelsel (omdat om stelsels te maken heb je een punt nodig
\newcommand{\pnR}{\rightarrow} %pijl naar Rechts
\newcommand{\limiet}{\lim\limits_}
\newcommand{\schrap}{\cancel}
\newcommand{\schrapb}{\bcancel} %backslash
\newcommand{\schrapk}{\xcancel} %kruis
\newcommand{\vervang}{\cancelto}
\newcommand{\intgrl}{\displaystyle\int} %integraal teken (bij het berekenen van primitieven)

\begin{document}

\noindent \large Auteur: Yoren Collier \\
\noindent \large Studierichting: Informatica\\
\noindent \large Oefeningenreeks 5A nr. 5.7\\

\medskip

\normalsize

$ BI \intgrl_\pi^{2\pi} \cos^2 \Lh x \Rh dx \adR OBI = \intgrl \cos^2 \Lh x \Rh dx$\\

$= \intgrl \dfrac{1 + \cos \Lh 2x \Rh}{2} dx = \dfrac{1}{2} \intgrl 1 + \cos \Lh 2x \Rh dx = \dfrac{1}{2} \Lv \intgrl1dx + \intgrl\cos \Lh 2x \Rh \Rv $\\

Stel: $t = 2x \adR dt = 2dx \asa dx = \dfrac{dt}{2}$\\

$= \dfrac{1}{2} \Lv x +c + \dfrac{1}{2} \intgrl \cos t dt \Rv = \dfrac{1}{2} \Lv x +c + \dfrac{1}{2} \sin 2x +c \Rv = \dfrac{1}{2}x +\dfrac{1}{4}\sin \Lh 2x \Rh +c$

$BI = \Lv \dfrac{1}{2}x +\dfrac{1}{4}\sin \Lh 2x \Rh \Rv_\pi^{2\pi} = \dfrac{1}{2} \Lh 2\pi \Rh + \dfrac{1}{4} \sin \Lh 2 \cdot 2\pi \Rh - \dfrac{1}{2} \Lh 2\pi \Rh - \dfrac{1}{4} \sin \Lh 2 \cdot 2\pi \Rh$\\

Doordat $\sin \Lh k \cdot 2\pi \Rh = 0$, met $k \in \mathbb{Z}$\\

$\dfrac{1}{2} \Lh \schrap{2}\pi - \schrap{\pi} +\schrapb{\dfrac{1}{2} \sin \Lh 4\pi \Rh} -\schrapb{\dfrac{1}{2} \sin \Lh 4\pi \Rh} \Rh = \dfrac{\pi}{2}$\\

\begin{thebibliography}{999}
%\bibitem{a} Referentie
\end{thebibliography}
\end{document} 